\documentclass[final]{beamer}

\usepackage[scale=1.2, size=a0, orientation=portrait]{beamerposter} % Use the beamerposter package for laying out the poster
\usepackage{epstopdf}
\usepackage{tipa}
\usepackage{booktabs}
\DeclareFontSubstitution{T3}{ptm}{m}{n}
\usetheme{confposter} % Use the confposter theme supplied with this template
\setbeamercolor{block title}{fg=ddblue,bg=white} % Colors of the block titles
\setbeamercolor{block body}{fg=black,bg=white} % Colors of the body of blocks
\setbeamercolor{block alerted title}{fg=white,bg=dblue!70} % Colors of the highlighted block titles
\setbeamercolor{block alerted body}{fg=black,bg=dblue!10} % Colors of the body of highlighted blocks
\setbeamertemplate{itemize items}[circle]
% Many more colors are available for use in beamerthemeconfposter.sty

%-----------------------------------------------------------
% Define the column widths and overall poster size
% To set effective sepwid, onecolwid and twocolwid values, first choose how many columns you want and how much separation you want between columns
% In this template, the separation width chosen is 0.024 of the paper width and a 4-column layout
% onecolwid should therefore be (1-(# of columns+1)*sepwid)/# of columns e.g. (1-(4+1)*0.024)/4 = 0.22
% Set twocolwid to be (2*onecolwid)+sepwid = 0.464
% Set threecolwid to be (3*onecolwid)+2*sepwid = 0.708

\newlength{\sepwid}
\newlength{\onecolwid}
\newlength{\twocolwid}
\newlength{\threecolwid}
\setlength{\paperwidth}{33.1in} % A0 width: 46.8in
\setlength{\paperheight}{46.8in} % A0 height: 33.1in
\setlength{\sepwid}{0.024\paperwidth} % Separation width (white space) between columns
\setlength{\onecolwid}{0.22\paperwidth} % Width of one column
\setlength{\twocolwid}{0.464\paperwidth} % Width of two columns
\setlength{\threecolwid}{0.708\paperwidth} % Width of three columns
\setlength{\topmargin}{-0.5in} % Reduce the top margin size
%-----------------------------------------------------------

\usepackage{graphicx}  % Required for including images

\usepackage{booktabs} % Top and bottom rules for tables

%-----------------------------
%	TITLE SECTION 
%-----------------------------

\title{Quantifying the information carried in tonal contrasts in Phom}

%-----------------------------

\setbeamertemplate{headline}{
 \leavevmode
  \begin{columns}
   \begin{column}{.2\linewidth}
    \centering
    \includegraphics[height=8cm]{EFLUlogo.jpg}
   \end{column}
   \begin{column}{.6\linewidth}
    \vskip1cm
    \centering
    \usebeamercolor{title in headline}{\color{jblue}\Huge{\inserttitle}\\[1ex]}

\begin{minipage}{0.5\textwidth}
  \centering
  \large{Phongshak Phom}\\[1ex]
  \normalsize{English and Foreign Languages University}\\[1ex]
  \normalsize{Hyderabad, TS, India}
\end{minipage}% 
\begin{minipage}{.5\textwidth}
  \centering
  \large{Charles Redmon}\\[1.5ex]
  \normalsize{University of Kansas}\\[1ex]
  \normalsize{Lawrence, KS, United States}
\end{minipage}
    \vskip1cm
    \large{The 22nd Himalayan Languages Symposium: IIT Guwahati, 2016}
    \vskip1cm
   \end{column}
   \begin{column}{.2\linewidth}
   \centering
   \includegraphics[height=6cm]{KUlogo.eps}
   \end{column}
   \vspace{1cm}
  \end{columns}
 \vspace{0.5in}
 \hspace{0.5in}\begin{beamercolorbox}[wd=32in,colsep=0.15cm]{cboxb}\end{beamercolorbox}
 \vspace{0.1in}
}

%------------------

\begin{document}

\addtobeamertemplate{block end}{}{\vspace*{2ex}} % White space under blocks
\addtobeamertemplate{block alerted end}{}{\vspace*{2ex}} % White space under highlighted (alert) blocks

\setlength{\belowcaptionskip}{2ex} % White space under figures
\setlength\belowdisplayshortskip{2ex} % White space under equations

\begin{frame}[t] % The whole poster is enclosed in one beamer frame

\begin{columns}[t] % The whole poster consists of three major columns, the second of which is split into two columns twice - the [t] option aligns each column's content to the top

%------------------------------------------------
%------------------------------------------------

\begin{column}{0.85\paperwidth} % Begin a column which is two columns wide (column 2)

\begin{columns}[t,totalwidth=0.85\paperwidth] % Split up the two columns wide column

\begin{column}{0.4\paperwidth}
\vspace{-.6in}


\begin{block}{Background}
  \begin{itemize}
    \item Phonemic contrasts are commonly treated as representing equivalent distinctions among abstract units in the \emph{inventory} of a language, but at their conception contrasts were fundamentally \emph{lexical} {\small (Martinet, 1938)}
    \item Within information theory, elements of the code are assumed to be utilized asymmetrically in conveying messages {\small (Shannon, 1948; Hockett, 1967)}
    \item Recent work has applied this perspective to cross-linguistic comparisons of contrast structures {\small (Surendran \& Niyogi, 2003; Oh et al., 2013, 2015)}, and to simulations of language change {\small (Wedel et al., 2013)}
    \item But the \emph{role} of contrasts is rarely studied in under-documented languages
    \item Phom is one such case: a Tibeto-Burman language spoken in Nagaland with a ternary (high, mid, low) lexical tone contrast {\small (Burling \& Phom, 1999)}
    \item Written Phom does not mark tones, making orthographic ambiguity one window on the unique information contributed by the tone system 
  \end{itemize}
\end{block}

%-----------------------------------
\end{column} 

\begin{column}{0.4\paperwidth}
\vspace{-.6in}

%-----------------------------------

\begin{block}{Methods}
  \begin{itemize}
    \item A 7,618-word corpus of written Phom based on selected chapters from \emph{Manshah} {\small (Phom, 2009)} was developed for this study
    \item 521 tonal minimal pairs were identified in the 2,635-word corpus-derived lexicon (all data were processed in Python 3.5 and analyzed in R 3.2)
    \item The contrast size $N_T$ -- the number of different lexical items (excluding homophones) represented by a given orthographic word -- was recorded for each \emph{token} ($n = 2,222$) of the 521 minimal pair \emph{types} in the corpus
    \item Values of $N_T$ were then updated as the first author was given the context in which these words occurred in the corpus in the following stepwise procedure:
    \begin{itemize}
      \item \emph{Unigram} (no context) $\rightarrow$ \emph{Bigram} (preceding word) $\rightarrow$ \emph{Trigram} (preceding 2 words)
    \end{itemize}
    \item Contrast size estimates defined a probability distribution from which effects of \textsc{Word Length} (syllable count, \emph{mono--tri}) and \textsc{Context} (N-gram size) on tonal disambiguation could be measured
  \end{itemize}
\end{block}

%-----------------------------------

\end{column} 

%-----------------------------------

\end{columns} % End of the split of column 2 - any content after this will now take up 2 columns width

%----------------------------------------
%	IMPORTANT RESULT
%----------------------------------------



%--------------------------------------------------
\begin{block}{Results}
\begin{columns}[t,totalwidth=0.85\paperwidth] % Split up the two columns wide column again

\begin{column}{0.4\paperwidth} % The first column within column 2 (column 2.1)

\begin{figure}
\includegraphics[width=0.7\linewidth]{plots/cdf.pdf}
\caption{CDFs of contrast size in unigram (black), bigram (orange), and trigram (blue) contexts.}
\end{figure}

\begin{itemize}
  \item Kolmogorov-Smirnov tests revealed a significant effect of \textsc{Context} on ambiguity: $trigram < bigram < unigram\, (p < 0.001)$
\end{itemize}
\vspace{1cm}

%-----------------------------------

\end{column} % End of column 2.1

\begin{column}{0.4\paperwidth} % The second column within column 2 (column 2.2)

%-----------------------------------
%	RESULTS
%-----------------------------------

\begin{itemize}
    \item Effects of \textsc{Word Length} were not consistent across contexts
    \begin{itemize}
      \item Unigram: $tri < di < mono$ in $N_T$ ($p_{m-d} = 0.019$; $p_{d-t} < 0.001$)
      \item Bigram: $tri < di = mono$ ($p_{d-t} = 0.006$; $p_{m-t} = 0.012$)
      \item Trigram: $tri < mono$ ($p_{m-t} = 0.025$)
    \end{itemize}
    \item Information gain ($\,N_{T(n-1)} - N_{T(n)}\,$) with context was also studied as a stochastic process, with gain at the bigram significantly greater than gain at the trigram ($p < 0.001$)
    \begin{itemize}
      \item This result was consistent across word lengths ($p < 0.001$)
      \item Trisyllables showed smaller gains relative to di/monosyllables at both bigram ($p < 0.05$) and trigram ($p < 0.01$) contexts, though this result is primarily due to their lower overall starting values of $N_T$
    \end{itemize}
    \item Context frequency in bigrams (i.e. frequency of the preceding word) was also shown to significantly modulate contrast size, with positive Kendall correlations between frequency and $N_T$ found both overall ($\tau = 0.10$, $p < 0.001$) and within word lengths ($0.06 < \tau < 0.14$, $p$s $< 0.01$)
    \item The inverse of contrast size ($P_T = 1/N_T$) was introduced as a measure of disambiguation probability (conceived as a lower bound on lexical predictability) for the study of tone system entropy
\end{itemize}


%------------------------------------

\end{column} % End of column 2.2

\end{columns} % End of the split of column 2
\end{block}

\begin{columns}[t,totalwidth=0.85\paperwidth] % Split up the two columns wide column

\begin{column}{0.4\paperwidth}
\vspace{-.6in}

\begin{block}{System Entropy}

\begin{figure}
\includegraphics[width=0.7\linewidth]{plots/KL.pdf}
\caption{Relative entropy (the Kullback-Leibler divergence, $P_{T(n-1)}||P_{T(n)}$) by word length.}
\end{figure}

\begin{itemize}
  \item Across word lengths, the Shannon entropy of $P_T$ decreased nonlinearly from 266 to 126 to 22 bits with increasing context (unigram--trigram)
  \item Trisyllables showed the greatest asymmetry in information gain with context, with the ratio of relative entropy from bigram to trigram in the following relation: $di\, (1.38) < mono\, (1.58) < tri\, (2.11)$
\end{itemize}

\end{block}

%-----------------------------------
\end{column} 

\begin{column}{0.4\paperwidth}
\vspace{-.6in}

%------------------------
%	SPECTRA
%------------------------
\begin{block}{Lexical Tone Distribution (monosyllables)}
\begin{figure}
\includegraphics[width=0.7\linewidth]{plots/tones.pdf}
\caption{Distribution of potential tonal variants of monosyllables among unique paths (i.e. distinct contrasts: HM, HL, HML, ML) from ambiguous unigram to disambiguated bi/tri-gram. }
\end{figure}

\begin{itemize}
  \item Binary contrasts involving the mid tone (HM, ML) were the least common, suggesting H and L might be more prominent under ambiguity
  \item Among disambiguated items, however, the three tones comprising the contrast were relatively evenly distributed: H -- 30\%; M -- 30\%; L -- 40\%
\end{itemize}

\bigskip
\bigskip

\end{block}

%-----------------------------------

\end{column} 

%-----------------------------------

\end{columns} % End of the split of column 2 - any content after this will now take up 2 columns width

\begin{columns}[t,totalwidth=0.85\paperwidth] % Split up the two columns wide column

\begin{column}{0.4\paperwidth}
\vspace{-.6in}

%------------------------
%	SPECTRA
%------------------------
\begin{block}{Morphology (disyllables)}
  \begin{itemize}
    \item Disyllabic words formed via compounding exhibit asymmetries in the tonal variant of each constituent
    \item For example, the mid-tone variant of the second constituent in \emph{yemshing} is more productive than the low (`stuffed up') or high (`to press') variants
  \end{itemize}

\bigskip
\bigskip

  \begin{tabular}{p{1in}lp{1in}lp{1in}l}
&    \emph{y\=em-sh\'ing} & & \emph{y\'em-sh\=\i ng} & & \emph{y\`em-sh\=\i ng}\\
&    search-press         & & dry-place        & & animal-place   \\
&    `to search for'      & & `a place to dry' & & `a place for animals' \\
  \end{tabular}

\bigskip
\bigskip

\begin{itemize}
  \item For \emph{vangdhum}, the first constituent is fixed, meaning the ambiguity is completely driven by the tone on the second constituent, \emph{dhum}
\end{itemize}

\bigskip
\bigskip

  \begin{tabular}{p{1in}lp{0.8in}l}
&    \emph{v\'ang-dh\`um} & & \emph{v\'ang-dh\'um} \\
&    rain-shade           & & rain-visit           \\
&    `shade, shelter'     & & `rain visitation' \\
  \end{tabular}

\end{block}

%-----------------------------------
\end{column} 

\begin{column}{0.4\paperwidth}
\vspace{-.6in}

%------------------------
%	SPECTRA
%------------------------
\begin{block}{Conclusions}
  \begin{itemize}
  \item The lexical role of the tone system in Phom, being more precisely quantified relative to effects of context (among others), may now serve as a reference for analyses of other tone systems in the region
    \item Future work should include syntactic and semantic constraints on $N_T$
  \end{itemize}
\end{block}

\begin{block}{References}
\footnotesize
\begin{center}
References will be made available upon request.
\end{center}
\end{block}

\begin{block}{Contact}
\footnotesize
  \begin{itemize}
    \item[E] Phongshak Phom: phongshak@gmail.com
    \item[] Charles Redmon: redmon@ku.edu
    \item[W] http://redmonc.github.io
  \end{itemize}
\end{block}


%-----------------------------------

\end{column} 

%-----------------------------------

\end{columns} % End of the split of column 2 - any content after this will now take up 2 columns width


\end{column} 



\end{columns} % End of all the columns in the poster

\end{frame} % End of the enclosing frame

\end{document}